\section{Evidence and Reading Discipline}

This chapter remains primarily a modelling chapter. Its formal claims are
proved only relative to:
\begin{enumerate}[nosep]
  \item the chapter's definitions;
  \item the monotone barrier scaffold from Chapter~\ref{ch:formal-language};
  \item the decision-load assumption introduced here.
\end{enumerate}

It does \emph{not} claim laboratory proof that every human plan reduces action
cost in every domain. The chapter proves what follows inside the model and no
more.

\section{Recommended Reading}

\begin{readingbox}
\textbf{Foundation}
\begin{itemize}[nosep]
  \item Steven Strogatz, \emph{Nonlinear Dynamics and Chaos}, for the logic of
        regime change and why small structural changes can matter more than raw
        force.
  \item Eugene Izhikevich, \emph{Dynamical Systems in Neuroscience}, for a more
        disciplined state-transition vocabulary than generic self-management
        language provides.
\end{itemize}
\textbf{Model discipline}
\begin{itemize}[nosep]
  \item Revisit Chapter~\ref{ch:formal-language} before reading this chapter as
        if it had established a calibrated behavioural law. It has not.
  \item Read Chapter~\ref{ch:dynamics} together with this chapter: planning is
        valuable because it changes later barrier structure before the decision
        window opens.
\end{itemize}
\textbf{Boundary}
\begin{itemize}[nosep]
  \item Empirical consciousness sources motivate taking structured state
        transitions seriously, but they are not direct proof of this chapter's
        planning propositions.
  \item Frontier physical theories are not needed to justify planning as an
        architectural control.
\end{itemize}
\end{readingbox}

\begin{summarybox}
This chapter now does less pretending and more actual work. It defines what a
plan changes, proves that resolving ambiguity and increasing preparation lower
future barrier under the model's monotone assumptions, proves a qualified
binary-execution advantage, and downgrades broader behavioural claims to the
level of disciplined interpretation. That is formally narrower than the
chapter's earlier rhetoric, but it is much closer to textbook honesty.
\end{summarybox}
